\documentclass[a4paper,11pt]{article}

\usepackage{revy}
\usepackage[utf8]{inputenc}
\usepackage[T1]{fontenc}
\usepackage[danish]{babel}
\usepackage{amsmath,amssymb}

\revyname{MatRevy}
\revyyear{2026}
\version{2026-10-23}
\eta{1 minutter}

\title{Pingu}
\author{Philip '23}

\begin{document}
\maketitle

\begin{roles}
\role{I}[Kisser] Instruktør
\role{P}[Daniel] Pingu
\role{F}[Niklas] Dum fysiker
\end{roles}

\begin{sketch}
\scene{Vi befinder os nede i sydenden. Pingu skal lave et epsilon-delta bevis. Pingu finder en tavle og struggler med at samle et kridt op i sine luffer. Det lykkedes endelig at få fat i kridtet og Pingu begynder at skrive definitionen på kontinuitet op. Fysikeren begynder på tavlen ved siden af at løse en opgave om heat equation hvor de betragter en pingvins varmetab til omgivelserne over tid. Pingu skæver interesseret over når det går op for ham at fysikerens opgave handler om pingviner. Han vender blikket tilbage til sin egen tavle. Fysikeren fjerner, for enkelthedens skyld, pingvinens luffer, og vælger til sidst at gøre den cylinderisk (Måske tegner de en cylinder på tavlen og skriver pingvin over). 
En ninja kommer forbi og fjerner Pingus luffer, så han ikke kan fortsætte med at lave sit bevis. Oprørt kigger han frem og tilbage mellem fysikerens tavle, sin krop, og hans Hagoromo kridt, som er faldet ned på gulvet:(
En ninja kommer bag Pingu, og løfter en mikrofon op til hans mund (fordi han ikke har arme længere). Pingu er nu, oprevet, sur, ked af det, og vender sig dramatisk mod publikum.}

\says{P}[Vredt, trist og dramatisk] NOOT NOOT!

\scene{Lys ned}
\end{sketch}
\end{document}
